\lamina{img/il_diapason.jpg}
\capitulo{CAPÍTULO 23}{EL ENTRELAZAMIENTO VERTICAL}{El diapasón invisible: sobre la fe, la oración y la meditación}{}
\noindent \capital{H}{ay} una palabra que este libro usa para cuando dos horizontes se acoplan: entrelazamiento. En el capítulo 12 hablamos de dos personas que llevan años entrelazadas sin quererlo. Ahora toca preguntar: ¿qué pasa si ese mismo enganche apunta hacia arriba, hacia algo que no está al otro lado de ninguna cama, sino más allá del propio borde?

Eso es, en el fondo, lo que hace rezar sin parar. Lo que hace meditar en serio. El ayuno, el rosario, el mantra, cantar en comunidad. En cualquier tradición que mires de cerca encontrarás la misma técnica: repetir hasta que el ruido se calla, quedarse quieto hasta que el sistema deja de inventar cosas nuevas, un ritmo compartido que sincroniza a toda una comunidad en la misma frecuencia.

\begin{fisica}
\lbl{En física esto se llama:} bajar la tasa de integración hasta casi cero.

\lbl{En la vida diaria es como:} afinar un instrumento con un diapasón que no puedes ver ni tocar. Ajustas la cuerda cada mañana sin saber si el diapasón sigue sonando.
\end{fisica}
Lo interesante no es la técnica. Es adónde apunta.

\seccion{Cuando se apunta al fondo}
\noindent Para el budismo y el hinduismo, se apunta directo al reservorio: esa plenitud sin forma de la que sale todo horizonte y a la que todo horizonte vuelve. El samadhi es el momento en que esa conexión se vuelve indistinguible de ser uno mismo con el fondo —la ola que descubre que es agua, no que viaja hacia el agua.

Aquí no hay Alguien al otro lado recibiendo el mensaje. Solo el descubrimiento de que nunca hubo dos lados.

\begin{fisica}
\lbl{En física esto se llama:} una correlación que se vuelve identidad.

\lbl{En la vida diaria es como:} la ola que descubre que es agua. No viaja hacia el mar: nunca salió de él.
\end{fisica}
\seccion{Cuando se apunta a Alguien}
\noindent Para el cristianismo y el islam, en cambio, no se apunta a ningún campo sin nombre. Se apunta a un horizonte anterior a todo lo demás, con voluntad propia, que sigue siendo Alguien por muy honda que llegue la oración. Rezar, aquí, no es fundirse. Es sostener una conexión con quien sigue estando al otro lado al final del proceso.

\begin{fisica}
\lbl{En física esto se llama:} un entrelazamiento donde ninguno de los dos se disuelve en el otro.

\lbl{En la vida diaria es como:} una carta que escribes cada noche a alguien que nunca contesta, y aun así sigues poniendo su nombre en el sobre.
\end{fisica}
Y quien no cree que haya nada más allá de lo físico hace la misma técnica, apuntando a nada más que a sí misma: entrelazamiento con el reservorio en el sentido más literal, porque para esa persona el reservorio es todo lo que hay del otro lado. Aquí no cabe la decepción, porque nunca hubo promesa de respuesta.

\seccion{Cinco cocinas, una receta}
\noindent Este libro no puede decir cuál de las cinco tiene razón. De un «es» nunca sale, por sí solo, un «debería» —eso ya lo vimos con Hume. Lo que sí puede decir es que cada tradición cierra esa brecha a su manera: el cristianismo con un vínculo —haces el bien porque ya te amaron, no para ganarte el amor—; el islam con una palabra revelada, tomada como un hecho; el hinduismo sin cerrarla desde fuera, porque el dharma no es una orden sino la descripción de cómo está tejido el mundo; el budismo mirando de frente el sufrimiento y su causa. Y quien no cree en nada trascendente la deja abierta, construyendo su propio código, sin que eso lo haga menos real para quien decide sostenerlo.

\begin{fisica}
\lbl{En software esto se llama:} cinco programas distintos que cumplen el mismo contrato.

\lbl{En la vida diaria es como:} cinco cocinas que sirven el mismo plato con cinco recetas. Desde la mesa, el sabor es idéntico. La diferencia está en la cocina.
\end{fisica}
Ninguna de esas cinco respuestas es la técnica en sí. La técnica —el silencio fabricado a mano, la repetición que vacía las palabras de peso— es igual en las cinco. Lo que las separa no se nota desde fuera, mirando a alguien de rodillas antes del amanecer. Se nota, si se nota, solo desde dentro. Y ni siquiera hay garantía de eso.

Queda una pregunta que este libro no puede resolver: si el entrelazamiento que termina en Alguien se siente distinto, desde dentro, del que termina en el agua sin nombre —o si la quietud es exactamente la misma quietud, y solo después, ya despiertos, decidimos cuál de las dos cosas acaba de pasar.

\begin{notacap}{Nota al capítulo 23}
\lbl{Lo que sí sabemos:} La práctica contemplativa aparece, con o sin nombre religioso, en todas las culturas humanas.

\lbl{Lo que no sabemos:} Si esa conexión hacia arriba corresponde a algo fuera del propio horizonte, sea con nombre o sin él.

\lbl{Preguntas que quedan:} ¿Se siente distinto, desde dentro, rezarle a Alguien que rezarle al agua? ¿O solo la interpretación posterior las separa?

\lbl{Si solo te quedas con una idea:} La técnica es la misma en las cinco tradiciones. Lo que cambia es a quién se imagina al otro lado del silencio.

\lbl{Lecturas:} Hume, D., \emph{Investigación sobre el entendimiento humano} (1748); Hesse, H., \emph{Siddhartha} (1922); Endō, S., \emph{Silencio} (1966).
\end{notacap}
